\silentchapter{7 Om det som ikkje lèt seg seie}{om det som ikkje lèt seg seie}

Om det som ikkje lèt seg seie, lyt formgjevaren forme.

Traktaten har skildra ei verd av former, krefter, landskap og agentar. Ho har gjeve ein formell likning for korleis form oppstår. Ho har definert omsorg som intelligensens innhald og intelligens som omsorgas form. Ho har vist at funksjonalismen var ein kollaps i omsorg, og at skaden ikkje let seg reparere.

Men traktaten har ikkje sagt noko om kvifor nokre former er vakre og andre ikkje er det. Ho har ikkje sagt noko om kvifor det å halde ein Wegner-stol i handa kjennest annleis enn å sitje i ein Monobloc. Ho har ikkje sagt noko om den stille tilfredsstillinga handverkaren kjenner når samanføyinga er presis, eller den naggande uroa designaren kjenner når noko ikkje stemmer men ho ikkje kan seie kva.

Dette er ikkje ein mangel ved rammeverket. Det er grensa hans.

Estetisk dom ligg utanfor formsystemet, på same vis som etikk ligg utanfor logikken. Ikkje fordi det er uviktig; det er kanskje det viktigaste. Men fordi det ikkje let seg fange i proposisjonar. Ein proposisjon kan seie at ein form okkuperer posisjon (52, 91, 48) i morforommet. Han kan seie at posisjonen ligg på ein haug i tilpassingslandskapet. Han kan seie at agentens lyskjegle romma fire dimensjonar då forma vart navigert. Men han kan ikkje seie at forma er god. Det siste er eit utsegn av ein heilt annan type.

Wittgenstein avslutta si eiga traktatform med å seie at det ein ikkje kan tale om, lyt ein teia om. For designaren er dette ikkje nok. Designarens svar på det useigelege er ikkje togn; det er handling. Ho formar. Ho svarar med hendene på det proposisjonane ikkje kan fange.

igskip

Rammeverket seier ikkje kva som er god form. Det seier ikkje kva ein bør designe. Det seier ikkje kva som er vakkert, kva som er rett, eller kva som er verdt å lage. Det seier ikkje noko om oppleving; om korleis det kjennest å sitje i ein stol, å gå gjennom ein bygning, å halde eit verktøy.

Det det seier, er dette: kvar form er ein posisjon i eit rom. Posisjonen er ikkje tilfeldig; ho er forma av krefter. Kreftene produserer eit landskap. Landskapet vert navigert av agentar med avgrensa lyskjegler. Agentens omsorg; mengda av dimensjonar ho held i merksemd; avgjer kor rik eller fattig forma vert.

Dersom rammeverket er rett, følgjer fleire ting for korleis designfaget bør organiserast:

Designutdanning bør lære landskapsnavigasjon, ikkje stilimitasjon. Studenten som forstår kreftene, kan navigere kvar som helst i landskapet. Studenten som berre kan imitere symptom, er bunden til dei posisjonane ho har lært.

Designhistorie bør studere seleksjonstrykk, ikkje stilperiodar. Stilperioden er eit symptom. Trykka er årsakene. Å organisere designhistoria etter stilperiodar er som å organisere medisinsk forsking etter symptom i staden for etter sjukdomsmekanismar.

Intelligens i design er omsorg: mengda av dimensjonar designaren held i hovudet samstundes. Denne mengda er trenbar. Lyskjegla kan utvidast. Kvar ny dimensjon som vert inkludert i kompromisset, produserer former som toler fleire trykk samstundes.

Skilet mellom industridesign, arkitektur, biologi og komputasjon er ein illusjon. Dei deler same formalisme: formrom, seleksjonstrykk, tilpassingslandskap, agentar med lyskjegler, grammatikkoperasjonar. Berre substratet er ulikt.

igskip

Denne boka er sjølv ein posisjon i eit formrom. Ho er eit forsøk på å kartleggje eit landskap som designteori har navigert i over hundre år utan å ha namna på kreftene. Kartet er uferdig. Somme regionar er godt lodda; dei empiriske testane i appendikset viser det. Andre er knapt skisserte. Modellen er konstruert for å kunne falsifiserast; dersom eitt einaste seleksjonstrykk viser seg å forklare all formvariasjon i ein klasse, kollapsar den logiske kjeda.

Men dersom modellen held; dersom det verkeleg er slik at form oppstår i feltet mellom agentar, navigert gjennom grammatikkoperasjonar i eit dynamisk landskap av seleksjonstrykk; så har me eit fundament. Ikkje ei samling av meiningane til framståande designarar. Ikkje ei samling av stilistiske preferansar. Eit fundament: ein struktur som kan berast av kven som helst, kvar som helst, uavhengig av substrat.

Rammeverket skal ikkje erstatte handverket. Det skal vise handverket kva det allereie gjer.

Om det som ikkje lèt seg seie, lyt formgjevaren forme.
